\documentclass[twoside,12pt]{book}

\usepackage[utf8]{inputenc}
\usepackage[T1]{fontenc}
\usepackage{mathpazo}
\usepackage{amsmath,amssymb,amsthm}
\usepackage{graphicx}
\usepackage{listings}
\usepackage{xcolor}
\usepackage{enumitem}
\usepackage{tikz}
\usepackage{algorithm}
\usepackage{algpseudocode}
\usepackage{booktabs}
\usepackage{geometry}
\usepackage{fancyhdr}
\usepackage{titlesec}
\usepackage[colorlinks=true, linkcolor=blue, citecolor=blue, urlcolor=blue, plainpages=false, pdfpagelabels]{hyperref}

\geometry{
  papersize={6in,9in},
  margin=0.75in,
  top=0.75in,bottom=0.75in,
  inner=0.85in,outer=0.65in,
  headheight=14pt
}

\pagestyle{fancy}
\fancyhf{}
\fancyhead[LE]{\small\textit{4. Moments and Deviations}}
\fancyhead[RO]{\small\textit{Randomized Algorithms}}
\fancyfoot[C]{\thepage}
\renewcommand{\headrulewidth}{0.4pt}

\definecolor{codegray}{rgb}{0.45,0.45,0.45}
\lstdefinestyle{cppstyle}{
  basicstyle=\ttfamily\scriptsize,
  keywordstyle=\bfseries,
  commentstyle=\itshape\color{codegray},
  numbers=left,
  numberstyle=\tiny\color{codegray},
  numbersep=8pt,
  frame=single,
  framerule=0.4pt,
  rulecolor=\color{codegray},
  backgroundcolor=\color{gray!5},
  breaklines=true,
  tabsize=4,
  showstringspaces=false,
  language=C++,
  morekeywords={nullptr,size_t,auto,const,constexpr,
    unordered_map,vector,pair,mt19937,
    uniform_int_distribution,INT_MAX}
}
\lstset{style=cppstyle}

\theoremstyle{plain}
\newtheorem{theorem}{Theorem}[chapter]
\newtheorem{lemma}[theorem]{Lemma}
\newtheorem{corollary}[theorem]{Corollary}
\newtheorem{proposition}[theorem]{Proposition}
\theoremstyle{definition}
\newtheorem{definition}[theorem]{Definition}
\newtheorem{example}[theorem]{Example}
\newtheorem{remark}[theorem]{Remark}
\theoremstyle{remark}
\newtheorem{exercise}[theorem]{Exercise}
\newtheorem{problem}[theorem]{Problem}
\newenvironment{cpplisting}{\begin{center}\small}{\end{center}}

\newcommand{\Exp}{\operatorname{\mathbb{E}}}
\newcommand{\Var}{\operatorname{Var}}
\newcommand{\Cov}{\operatorname{Cov}}

\begin{document}

\chapter*{4 Moments and Deviations}
\addcontentsline{toc}{chapter}{4 Moments and Deviations}
\markboth{4 Moments and Deviations}{Randomized Algorithms}
\label{ch:moments}

In the preceding chapters we studied some simple randomized algorithms and encountered the need to analyze the probability that a random variable deviates significantly from its expected value. We now develop general tools for bounding tail probabilities---the probability that a random variable deviates by a given amount from its expectation. The probability of such a deviation is referred to as a \emph{tail probability}. Readers wishing to review basic material on probability and distributions may consult Appendix~A.

\section{Occupancy Problems}
\label{sec:occupancy}

We begin with an example of an occupancy problem. In such problems we envision each of $m$ indistinguishable objects (``balls'') being randomly assigned to one of $n$ distinct classes (``bins''). In other words, each ball is placed in a bin chosen independently and uniformly at random. We are interested in questions such as: what is the maximum number of balls in any bin? what is the expected number of bins with $k$ balls in them? Such problems are at the core of the analyses of many randomized algorithms ranging from data structures to routing in parallel computers.

Our discussion of the occupancy problem will illustrate a recurrent tool in the analysis of randomized algorithms: that the probability of the union of events is no more than the sum of their probabilities. This is a special case of the Boole--Bonferroni Inequalities (Proposition~C.2) and can be formally stated as follows: for arbitrary events $\mathcal{E}_1, \mathcal{E}_2, \ldots, \mathcal{E}_n$, not necessarily independent,
\[
\Pr\Bigl[\bigcup_{i=1}^{n} \mathcal{E}_i\Bigr] \leq \sum_{i=1}^{n} \Pr[\mathcal{E}_i].
\]
This principle is extremely useful because it assumes nothing about the dependencies between the events.

Consider first the case $m = n$. For $1 \leq i \leq n$, let $X_i$ be the number of balls in the $i$th bin. Following Example~1.1, we have $\mathbb{E}[X_i] = 1$ for all $i$. Yet we do not expect that during a typical experiment every bin receives exactly one ball. Rather, we expect some bins to have no balls at all, and others to have many more than one.

Let us try now to make a statement of the form ``with very high probability, no bin receives more than $k$ balls,'' for a suitably chosen $k$. Let $\mathcal{E}_j(k)$ denote the event that bin $j$ has $k$ or more balls in it. We concentrate on analyzing $\mathcal{E}_1(k)$.

The probability that bin~1 receives exactly $i$ balls is
\[
\Pr[X_1 = i] = \binom{m}{i} \frac{1}{n^i}\Bigl(1 - \frac{1}{n}\Bigr)^{m-i} \leq \frac{m^i}{i! \cdot n^i}.
\]
The second inequality results from an upper bound for binomial coefficients (Proposition~B.2). Thus,
\[
\Pr[\mathcal{E}_1(k)] \leq \sum_{i=k}^{m} \frac{m^i}{i! \cdot n^i} \leq \frac{(m/n)^k}{k!} \cdot \frac{1}{1 - m/(kn)}.
\]
Setting $k^* = \lceil (en \ln n) / \ln \ln n \rceil$, we obtain
\begin{equation}
\Pr[\mathcal{E}_1(k^*)] \leq \frac{(e \ln n)^{k^*}}{(k^*)!} \leq \frac{1}{n^2}.
\label{eq:occ-bound}
\end{equation}

The same computation tells us that this upper bound applies to $\Pr[\mathcal{E}_i(k^*)]$ for all~$i$, but can we say that no bin is likely to have more than $k^*$ balls in it? For this we invoke the principle mentioned at the beginning of this section: the probability of the union of the events $\mathcal{E}_i(k^*)$ is no more than their sum.

\begin{theorem}[Maximum Occupancy]
\label{thm:occ-max}
With probability at least $1 - 1/n$, no bin has more than $k^* = \lceil (en \ln n) / \ln \ln n \rceil$ balls when $m = n$ balls are randomly assigned to $n$ bins.
\end{theorem}

\begin{proof}
By the union bound and~\eqref{eq:occ-bound},
\[
\Pr\Bigl[\bigcup_{i=1}^{n} \mathcal{E}_i(k^*)\Bigr] \leq \sum_{i=1}^{n} \Pr[\mathcal{E}_i(k^*)] \leq n \cdot \frac{1}{n^2} = \frac{1}{n}.
\]
\end{proof}

\medskip

Interestingly, when $m$ is of the order of $n \ln n$, the bin with the most balls has about the same number of balls as the expected number of balls in any bin. This phenomenon is exploited in a number of randomized algorithms.

\begin{exercise}
For $m = n \log n$, show that with probability $1 - o(1)$ every bin contains $O(\log n)$ balls.
\end{exercise}

\subsection{The Birthday Problem}

We turn to a classic combinatorial problem. Suppose that $m$ balls are randomly assigned to $n$ bins. We study the probability of the event that they all land in distinct bins. The special case $n = 365$ is popular in mathematical lore as the \emph{birthday problem}. The interpretation is that the $365$ days of the year correspond to $365$ bins, and the birthday of each of $m$ people is chosen independently and uniformly from all $365$ days (ignoring leap years). How large must $m$ be before two people in the group are likely to share their birthdays?

Consider the assignment of the balls to the bins as a sequential process: we throw the first ball into a random bin, then the second ball, and so on. For $2 \leq i \leq m$, let $\mathcal{E}_i$ denote the event that the $i$th ball lands in a bin not containing any of the first $i-1$ balls. We will bound $\Pr[\bigcap_{i=2}^{m} \mathcal{E}_i]$ from above. From the definition of conditional probability,
\[
\Pr\Bigl[\bigcap_{i=2}^{m} \mathcal{E}_i\Bigr] = \Pr[\mathcal{E}_2] \Pr[\mathcal{E}_3 \mid \mathcal{E}_2] \Pr[\mathcal{E}_4 \mid \mathcal{E}_2 \cap \mathcal{E}_3] \cdots \Pr[\mathcal{E}_m \mid \bigcap_{i=2}^{m-1} \mathcal{E}_i].
\]
Now, $\Pr[\mathcal{E}_i \mid \bigcap_{j=2}^{i-1} \mathcal{E}_j] = 1 - (i-1)/n$. Making use of the fact that $1 - x \leq e^{-x}$, we have
\[
\Pr\Bigl[\bigcap_{i=2}^{m} \mathcal{E}_i\Bigr] \leq \prod_{i=2}^{m}\Bigl(1 - \frac{i-1}{n}\Bigr) \leq \prod_{i=2}^{m} e^{-(i-1)/n} = e^{-\sum_{i=1}^{m-1} i/n} = e^{-m(m-1)/(2n)}.
\]
Thus, we see that for $m$ equal to $\lfloor \sqrt{2n+1} \rfloor$, the probability that all $m$ balls land in distinct bins is at most $1/e$; as $m$ increases beyond this value, the probability drops rapidly.

\begin{cpplisting}
\label{code:birthday}
\begin{lstlisting}[language=C++, frame=single, framerule=0.4pt,
  rulecolor=\color{gray}, backgroundcolor=\color{gray!5},
  basicstyle=\ttfamily\footnotesize, keywordstyle=\bfseries,
  commentstyle=\itshape\color{gray}, numbers=left,
  numberstyle=\tiny\color{gray}, numbersep=8pt,
  tabsize=2, showstringspaces=false]
#include <iostream>
#include <vector>
#include <random>
#include <cmath>

// Estimate the probability of a birthday collision among m people
// with n possible birthdays, using repeated simulations.
double birthday_collision_prob(int m, int n, int trials,
                               std::mt19937& rng) {
    int collisions = 0;
    std::uniform_int_distribution<int> dist(0, n - 1);
    for (int t = 0; t < trials; ++t) {
        std::vector<bool> seen(n, false);
        bool collision = false;
        for (int i = 0; i < m; ++i) {
            int bday = dist(rng);
            if (seen[bday]) { collision = true; break; }
            seen[bday] = true;
        }
        if (collision) ++collisions;
    }
    return static_cast<double>(collisions) / trials;
}

int main() {
    std::mt19937 rng(42);
    int n = 365;
    // Theoretical threshold: m ~ sqrt(2 * n * ln(1/(1-p)))
    // For p = 0.5, m ~ 22.49
    double analytical = 1.0 - std::exp(-0.5);
    std::cout << "Birthday Problem (n=" << n << "):\n";
    for (int m : {10, 15, 20, 23, 30, 50}) {
        double sim = birthday_collision_prob(m, n, 100000, rng);
        double theory = 1.0 - std::exp(-static_cast<double>(m * (m - 1)) / (2.0 * n));
        std::cout << "  m=" << m
                  << "  simulated=" << sim
                  << "  theoretical=" << theory << "\n";
    }
    return 0;
}
\end{lstlisting}
\end{cpplisting}

\section{The Markov and Chebyshev Inequalities}
\label{sec:markov-chebyshev}

We have seen above that making statements about the probability that a random variable deviates far from its expectation may involve a detailed, problem-specific analysis. Often, one can avoid such detailed analyses by resorting to general inequalities on such tail probabilities.

We begin with the Markov inequality, a fundamental tool we will invoke repeatedly when we develop more sophisticated bounding techniques.

\begin{theorem}[Markov Inequality]
\label{thm:markov}
Let $Y$ be a random variable assuming only non-negative values. Then for all $t \in \mathbb{R}^+$,
\[
\Pr[Y \geq t] \leq \frac{\mathbb{E}[Y]}{t}.
\]
Equivalently, $\Pr[Y \geq k\,\mathbb{E}[Y]] \leq 1/k$.
\end{theorem}

\begin{proof}
Define a function $f(y)$ by $f(y) = 1$ if $y \geq t$, and $0$ otherwise. Then $\Pr[Y \geq t] = \mathbb{E}[f(Y)]$. Since $f(y) \leq y/t$ for all $y$,
\[
\mathbb{E}[f(Y)] \leq \mathbb{E}\Bigl[\frac{Y}{t}\Bigr] = \frac{\mathbb{E}[Y]}{t},
\]
and the theorem follows.
\end{proof}

This is the tightest possible bound when we know only that $Y$ is non-negative and has a given expectation. Unfortunately, the Markov inequality by itself is often too weak to yield useful results.

\begin{exercise}
Given a positive integer $k$, describe a random variable $X$ assuming only non-negative values, such that $\Pr[X \geq k\,\mathbb{E}[X]] = 1/k$.
\end{exercise}

\begin{exercise}
Let $Y$ be any random variable and $h$ any non-negative real function. Show that for all $t \in \mathbb{R}^+$,
\[
\Pr[h(Y) \geq t] \leq \frac{\mathbb{E}[h(Y)]}{t}.
\]
\end{exercise}

The following generalization of Markov's inequality underlies its usefulness in deriving stronger bounds.

We now show that the Markov inequality can be used to derive better bounds on the tail probability by using more information about the distribution of the random variable. The first of these is the Chebyshev bound, which is based on the knowledge of the variance of the distribution.

For a random variable $X$ with expectation $\mu_X$, its \emph{variance} $\sigma_X^2$ is defined to be $\mathbb{E}[(X - \mu_X)^2]$. The \emph{standard deviation} of $X$, denoted $\sigma_X$, is the positive square root of $\sigma_X^2$.

\begin{theorem}[Chebyshev's Inequality]
\label{thm:chebyshev}
Let $X$ be a random variable with expectation $\mu_X$ and standard deviation $\sigma_X$. Then for any $t \in \mathbb{R}^+$,
\[
\Pr[|X - \mu_X| \geq t\,\sigma_X] \leq \frac{1}{t^2}.
\]
\end{theorem}

\begin{proof}
First, note that
\[
\Pr[|X - \mu_X| \geq t\,\sigma_X] = \Pr[(X - \mu_X)^2 \geq t^2 \sigma_X^2].
\]
The random variable $Y = (X - \mu_X)^2$ has expectation $\sigma_X^2$, and applying the Markov inequality to $Y$ bounds this probability from above by $1/t^2$.
\end{proof}

\begin{exercise}
The use of the variance of a random variable in bounding its deviation from its expectation is called the \emph{second moment method}. In an analogous fashion, we can speak of the \emph{$k$th moment method}: let $k$ be even, and suppose we have a random variable $X$ for which $\mu_k = \mathbb{E}[(X - \mu_X)^k]$ exists. Show that
\[
\Pr[|X - \mu_X| > t\,(\mu_k)^{1/k}] \leq \frac{1}{t^k}.
\]
Why is the $k$th moment method difficult to invoke for odd values of $k$?
\end{exercise}

\section{Randomized Selection}
\label{sec:rand-sel}

We now consider the use of random sampling for the problem of selecting the $k$th smallest element in a set $S$ of $n$ elements drawn from a totally ordered universe. We assume that the elements of $S$ are all distinct, although it is not very hard to modify the following analysis to allow for multisets. Let $r_S(t)$ denote the rank of an element $t$ (the $k$th smallest element has rank~$k$) and let $S_{\langle i \rangle}$ denote the $i$th smallest element of $S$.

\begin{algorithm}[H]
\caption{LazySelect}
\label{alg:lazy-select}
\begin{algorithmic}[1]
\Require A set $S$ of $n$ elements from a totally ordered universe; integer $k \in [n]$
\Ensure The $k$th smallest element of $S$, $S_{\langle k \rangle}$
\State Pick $n^{3/4}$ elements from $S$, chosen independently and uniformly at random with replacement; call this multiset of elements $R$
\State Sort $R$ using any optimal sorting algorithm in $O(n^{3/4} \log n)$ steps
\State Let $x = k \cdot n^{-1/4}$. Set $t = \max\{\lfloor x - \sqrt{n} \rfloor, 1\}$ and $h = \min\{\lceil x + \sqrt{n} \rceil, n^{3/4}\}$
\State Let $a = R_{\langle t \rangle}$ and $b = R_{\langle h \rangle}$. By comparing $a$ and $b$ to every element of $S$, determine $r_S(a)$ and $r_S(b)$
\If{$k < n^{1/4}$}
    \State $P \gets \{y \in S \mid y \leq b\}$
\ElsIf{$k > n - n^{1/4}$}
    \State $P \gets \{y \in S \mid y \geq a\}$
\Else
    \State $P \gets \{y \in S \mid a \leq y \leq b\}$
\EndIf
\State Check whether $S_{\langle k \rangle} \in P$ and $|P| \leq 4n^{3/4} + 2$. If not, repeat Steps~1--5 until such a set $P$ is found
\State Sort $P$ and identify $P_{\langle k - r_S(a) + 1 \rangle}$, which is $S_{\langle k \rangle}$
\end{algorithmic}
\end{algorithm}

The idea of the algorithm is to identify two elements $a$ and $b$ in $S$ such that both of the following hold with high probability:
\begin{enumerate}
\item The element $S_{\langle k \rangle}$ that we seek is in $P$.
\item The set $P$ of elements between $a$ and $b$ is not very large, so that we can sort $P$ inexpensively in Step~7.
\end{enumerate}

We examine how either of these requirements could fail. We focus on the most interesting case when $k \in [n^{1/4}, n - n^{1/4}]$, so that $P = \{y \in S \mid a \leq y \leq b\}$.

If the element $a$ is greater than $S_{\langle k \rangle}$ (or if $b$ is smaller than or equal to $S_{\langle k \rangle}$), we fail because $P$ does not contain $S_{\langle k \rangle}$. For this to happen, fewer than $t$ of the samples in $R$ should be smaller than $S_{\langle k \rangle}$ (respectively, at least $h$ of the random samples should be smaller than $S_{\langle k \rangle}$). We will bound the probability that this happens using the Chebyshev bound.

\begin{lemma}[Additivity of Variances]
\label{lemma:var-sum}
Let $X_1, X_2, \ldots, X_m$ be independent random variables. Let $X = \sum_{i=1}^{m} X_i$. Then $\sigma_X^2 = \sum_{i=1}^{m} \sigma_{X_i}^2$.
\end{lemma}

\begin{proof}
Let $\mu_i$ denote $\mathbb{E}[X_i]$, and $\mu = \sum_{i=1}^{m} \mu_i$. The variance of $X$ is given by
\[
\mathbb{E}[(X - \mu)^2] = \mathbb{E}\Bigl[\Bigl(\sum_{i=1}^{m} (X_i - \mu_i)\Bigr)^2\Bigr].
\]
Expanding the latter and using linearity of expectations,
\[
\mathbb{E}[(X - \mu)^2] = \sum_{i=1}^{m} \mathbb{E}[(X_i - \mu_i)^2] + 2 \sum_{i < j} \mathbb{E}[(X_i - \mu_i)(X_j - \mu_j)].
\]
Since all pairs $X_i, X_j$ are independent, so are the pairs $(X_i - \mu_i), (X_j - \mu_j)$. By the fact that expectation factors for independent variables, each term in the latter summation equals $\mathbb{E}[(X_i - \mu_i)] \cdot \mathbb{E}[(X_j - \mu_j)]$. Since $\mathbb{E}[(X_i - \mu_i)] = 0$, the latter summation vanishes.
\end{proof}

\begin{theorem}[LazySelect Correctness]
\label{thm:lazy-select}
With probability $1 - O(n^{-1/4})$, LazySelect finds $S_{\langle k \rangle}$ on the first pass through Steps~1--7, and thus performs only $2n + o(n)$ comparisons.
\end{theorem}

\begin{proof}
Step~3 requires $2n$ comparisons, and all other steps perform $o(n)$ comparisons, provided the algorithm finds $S_{\langle k \rangle}$ on the first pass. We now consider the first mode of failure: $a > S_{\langle k \rangle}$ because fewer than $t$ of the samples in $R$ are less than or equal to $S_{\langle k \rangle}$.

Let $X_i = 1$ if the $i$th random sample is at most $S_{\langle k \rangle}$, and $0$ otherwise; thus $\Pr[X_i = 1] = k/n$, and $\Pr[X_i = 0] = 1 - k/n$. Let $X = \sum_{i=1}^{n^{3/4}} X_i$ be the number of samples of $R$ that are at most $S_{\langle k \rangle}$. The random variables $X_i$ are Bernoulli trials. Using Lemma~\ref{lemma:var-sum} and the variance of a Bernoulli trial with success probability $p$,
\[
\mu_X = k n^{3/4} / n = k n^{-1/4},
\]
and
\[
\sigma_X^2 = n^{3/4} \cdot \frac{k}{n} \cdot \Bigl(1 - \frac{k}{n}\Bigr) \leq n^{3/4} \cdot \frac{1}{4}.
\]
This implies that $\sigma_X \leq n^{3/8}/2$. Applying the Chebyshev bound to $X$,
\[
\Pr[|X - \mu_X| \geq \sqrt{n}] = \Pr[|X - \mu_X| \geq 2n^{1/8} \sigma_X] = o(n^{-1/4}).
\]
An essentially identical argument shows that $\Pr[b < S_{\langle k \rangle}] = o(n^{-1/4})$. Since the probability of the union of events is at most the sum of their probabilities, the probability that either of these events occurs (causing $S_{\langle k \rangle}$ to lie outside $P$) is $O(n^{-1/4})$.

The analysis of the second mode of failure---that $P$ contains more than $4n^{3/4} + 2$ elements---is very similar, with $k_t$ and $k_h$ playing the role of $k$. Adding up the probabilities of all failure modes, we find that the probability that Steps~1--5 fail to find a suitable set $P$ is $O(n^{-1/4})$.
\end{proof}

This adds to the significance of LazySelect: the best known deterministic selection algorithms use $3n$ comparisons in the worst case and are quite complicated to implement. Further, it is known that any deterministic algorithm for finding the median requires at least $2n$ comparisons, so we have a randomized algorithm that is both fast and has an expected number of comparisons that is provably smaller than that of any deterministic algorithm.

\begin{exercise}
The failure probability can be driven down further at the expense of increased running time. For a suitable definition of the $o(n)$ term, give an upper bound on the probability that the algorithm does not find $S_{\langle k \rangle}$ in $cn + o(n)$ steps for $c > 2$.
\end{exercise}

\begin{exercise}
Suggest a modification in the algorithm that brings the constant in the linear term down to $1.5$ from $2$. We will refine this further in Problem~4.16.
\end{exercise}

\begin{exercise}
Show that as a direct corollary of Theorem~\ref{thm:lazy-select}, the expected running time of the LazySelect algorithm is $2n + o(n)$.
\end{exercise}

\begin{cpplisting}
\label{code:lazy-select}
\begin{lstlisting}[language=C++, frame=single, framerule=0.4pt,
  rulecolor=\color{gray}, backgroundcolor=\color{gray!5},
  basicstyle=\ttfamily\footnotesize, keywordstyle=\bfseries,
  commentstyle=\itshape\color{gray}, numbers=left,
  numberstyle=\tiny\color{gray}, numbersep=8pt,
  tabsize=2, showstringspaces=false]
#include <algorithm>
#include <cmath>
#include <iostream>
#include <random>
#include <vector>

// Returns the k-th smallest element (0-indexed) using LazySelect.
// Expected comparisons: 2n + o(n), succeeding w.h.p. on first pass.
int lazy_select(std::vector<int>& S, int k, std::mt19937& rng) {
    int n = static_cast<int>(S.size());
    int sample_size = static_cast<int>(std::ceil(std::pow(n, 0.75)));
    std::uniform_int_distribution<int> dist(0, n - 1);

    while (true) {
        // Step 1: sample n^{3/4} elements with replacement
        std::vector<int> R(sample_size);
        for (int i = 0; i < sample_size; ++i)
            R[i] = S[dist(rng)];

        // Step 2: sort R
        std::sort(R.begin(), R.end());

        // Step 3: compute indices t, h
        double x = static_cast<double>(k) * std::pow(n, -0.25);
        int t = std::max(static_cast<int>(std::floor(x - std::sqrt(n))), 1);
        int h = std::min(static_cast<int>(std::ceil(x + std::sqrt(n))),
                         sample_size - 1);
        int a_val = R[t];
        int b_val = R[h];

        // Step 4: compute ranks of a and b in S
        int rank_a = 0, rank_b = 0;
        for (int i = 0; i < n; ++i) {
            if (S[i] <= a_val) ++rank_a;
            if (S[i] < b_val) ++rank_b;
        }

        // Build candidate set P
        std::vector<int> P;
        if (k < static_cast<int>(std::pow(n, 0.25))) {
            for (int v : S) if (v <= b_val) P.push_back(v);
        } else if (k > n - static_cast<int>(std::pow(n, 0.25))) {
            for (int v : S) if (v >= a_val) P.push_back(v);
        } else {
            for (int v : S)
                if (v >= a_val && v <= b_val) P.push_back(v);
        }

        // Check S<k> is in P and |P| is small enough
        int target_in_P = k - rank_a + 1;
        if (target_in_P >= 1 && target_in_P <= static_cast<int>(P.size())
            && static_cast<int>(P.size()) <= 4 * sample_size + 2) {
            // Step 7: sort P and pick the element
            std::sort(P.begin(), P.end());
            return P[target_in_P - 1];
        }
        // Otherwise retry
    }
}

int main() {
    std::mt19937 rng(42);
    std::vector<int> S(10000);
    for (int i = 0; i < 10000; ++i) S[i] = i + 1;
    std::shuffle(S.begin(), S.end(), rng);

    int k = 5000;  // find the median
    int result = lazy_select(S, k, rng);
    std::cout << "LazySelect: " << k << "th smallest = "
              << result << "\n";

    std::sort(S.begin(), S.end());
    std::cout << "Sorted:     " << k << "th smallest = "
              << S[k - 1] << "\n";
    return 0;
}
\end{lstlisting}
\end{cpplisting}

\section{Two-Point Sampling}
\label{sec:two-point}

We have so far been making use of the fact that the variance of the sum of independent random variables equals the sum of their variances. In fact, we can make a stronger statement. Let $X$ and $Y$ be discrete random variables defined on the same probability space. The \emph{joint density function} of $X$ and $Y$ is the function
\[
p(x, y) = \Pr[\{X = x\} \cap \{Y = y\}].
\]
A set of random variables $X_1, X_2, \ldots, X_m$ is said to be \emph{pairwise independent} if for all $i \neq j$, and $x, y \in \mathbb{R}$,
\[
\Pr[X_i = x \mid X_j = y] = \Pr[X_i = x].
\]

\begin{exercise}
Let $n$ be a prime number and $\mathbb{Z}_n$ denote the field of integers modulo~$n$. For $a$ and $b$ chosen independently and uniformly at random from $\mathbb{Z}_n$, let $Y_i = ai + b \pmod{n}$. Show that for $i \neq j \pmod{n}$, $Y_i$ and $Y_j$ are uniformly distributed on $\mathbb{Z}_n$ and pairwise independent. (Make use of the fact that in the field $\mathbb{Z}_n$, given fixed values for $Y_i$ and $Y_j$, we can solve $Y_i = ai + b \pmod{n}$ and $Y_j = aj + b \pmod{n}$ uniquely for $a$ and~$b$.)
\end{exercise}

\begin{exercise}
Let $X_1, X_2, \ldots, X_m$ be pairwise independent random variables, and $X = \sum_{i=1}^{m} X_i$. Show that $\sigma_X^2 = \sum_{i=1}^{m} \sigma_{X_i}^2$.
\end{exercise}

We now consider an application of these concepts to the reduction of the random bits used by RP algorithms. Consider an RP algorithm $A$ for deciding whether input strings $x$ belong to a language~$L$. Given $x$, $A$ picks a random number $r$ from the range $\mathbb{Z}_n = \{0, \ldots, n-1\}$ and computes a binary value $A(x, r)$ with the following properties:
\begin{itemize}
\item If $x \in L$, then $A(x, r) = 1$ for at least half the possible values of $r$.
\item If $x \notin L$, then $A(x, r) = 0$ for all possible choices of $r$.
\end{itemize}

For any $x \in L$, we refer to the values of $r$ for which $A(x, r) = 1$ as \emph{witnesses} for $x$; clearly, at least $n/2$ of the $n$ possible values of $r$ are witnesses. Choosing $t$ independent random numbers is expensive in that it requires $O(t \log n)$ random bits. Suppose instead that we are only willing to use $O(\log n)$ random bits---in particular, suppose we wish to use only two independent samples from $\mathbb{Z}_n$.

The naive usage of $a$ and $b$ as potential witnesses yields an upper bound of only $1/4$ on the probability of incorrect classification. Here is a better scheme: let $r_i = ai + b \pmod{n}$, and compute $A(x, r_i)$ for $1 \leq i \leq t$. As before, if for any $i$ we obtain $A(x, r_i) = 1$, we declare that $x$ is in $L$, else we declare that $x$ is not in $L$.

\begin{theorem}[Probability Amplification via Two-Point Sampling]
\label{thm:two-point}
With the two-point sampling scheme, the error probability is at most $1/t$.
\end{theorem}

\begin{proof}
We need to worry about the possibility of making error only when the input $x$ is in $L$. By the result of Exercise~3.7, the random $r_i$'s are pairwise independent and, therefore, so are the random variables $A(x, r_i)$, for $1 \leq i \leq t$. Let $Y = \sum_{i=1}^{t} A(x, r_i)$. Assuming that $x \in L$, $\mathbb{E}[Y] \geq t/2$ and $\sigma_Y^2 \leq t/4$, so $\sigma_Y \leq \sqrt{t}/2$. The probability that the pairwise independent iterations produce an incorrect classification corresponds to the event $\{Y = 0\}$, and
\[
\Pr[Y = 0] \leq \Pr[|Y - \mathbb{E}[Y]| \geq t/2].
\]
By the Chebyshev inequality, the latter is at most $1/t$. Thus, the error probability is at most $1/t$, which is a considerable improvement over the error bound of $1/4$ achieved by the naive use of $a$ and~$b$.
\end{proof}

\begin{cpplisting}
\label{code:two-point}
\begin{lstlisting}[language=C++, frame=single, framerule=0.4pt,
  rulecolor=\color{gray}, backgroundcolor=\color{gray!5},
  basicstyle=\ttfamily\footnotesize, keywordstyle=\bfseries,
  commentstyle=\itshape\color{gray}, numbers=left,
  numberstyle=\tiny\color{gray}, numbersep=8pt,
  tabsize=2, showstringspaces=false]
#include <iostream>
#include <random>
#include <vector>

// Simulated RP algorithm: for x in L, A(x,r)=1 for half the r's.
// For x not in L, A(x,r)=0 for all r.
bool RP_algorithm(int x, int r, bool x_in_L) {
    if (!x_in_L) return false;
    // Deterministic witness check: r < n/2 means witness
    return r < 500;
}

// Naive: use a and b directly as witnesses. Error <= 1/4.
bool naive_two_point(int x, int n, bool x_in_L,
                     std::mt19937& rng) {
    std::uniform_int_distribution<int> dist(0, n - 1);
    int a = dist(rng), b = dist(rng);
    if (RP_algorithm(x, a, x_in_L)) return true;
    if (RP_algorithm(x, b, x_in_L)) return true;
    return false;
}

// Amplified: use r_i = a*i + b mod n for t samples.
// Error <= 1/t by pairwise independence.
bool amplified_two_point(int x, int n, int t, bool x_in_L,
                         std::mt19937& rng) {
    std::uniform_int_distribution<int> dist(0, n - 1);
    int a = dist(rng), b = dist(rng);
    for (int i = 0; i < t; ++i) {
        int ri = (static_cast<long long>(a) * i + b) % n;
        if (RP_algorithm(x, ri, x_in_L)) return true;
    }
    return false;
}

int main() {
    std::mt19937 rng(42);
    int n = 1000;
    int trials = 100000;
    int t = 10;  // amplification parameter

    // Test with x in L (should almost always return true)
    int naive_false = 0, amp_false = 0;
    for (int i = 0; i < trials; ++i) {
        if (!naive_two_point(1, n, true, rng)) ++naive_false;
        if (!amplified_two_point(1, n, t, true, rng)) ++amp_false;
    }
    std::cout << "x in L:\n"
              << "  Naive error rate:   "
              << static_cast<double>(naive_false) / trials << "\n"
              << "  Amplified error rate: "
              << static_cast<double>(amp_false) / trials << "\n"
              << "  (theory: naive <= 1/4, amplified <= 1/t = "
              << 1.0 / t << ")\n";
    return 0;
}
\end{lstlisting}
\end{cpplisting}

\section{The Stable Marriage Problem}
\label{sec:stable-marriage}

Consider a society in which there are $n$ men (denoted by capital letters $A, B, C, \ldots$) and $n$ women (denoted by lowercase letters $a, b, c, \ldots$). A marriage $M$ is a one-to-one correspondence between the men and the women. Assume a monogamous society. Each person has a preference list of the members of the opposite sex organized in decreasing order of desirability. A marriage is said to be \emph{unstable} if there exist two married couples $X$--$x$ and $Y$--$y$ such that $X$ desires $y$ more than $x$, and $y$ desires $X$ more than $Y$, implying that $X$--$y$ will have a tendency to leave their current mates to marry each other. A marriage $M$ in which there are no dissatisfied couples is called a \emph{stable marriage}.

\begin{example}
For $n = 4$, consider the following preference lists.
\[
\begin{array}{llll}
A: a\; b\; c\; d & B: b\; a\; e\; d & C: a\; d\; e\; b & D: d\; e\; a\; b \\
a: A\; B\; C\; D & b: D\; C\; B\; A & c: A\; B\; C\; D & d: C\; D\; A\; B
\end{array}
\]
Consider the marriage $M$ given by $A$--$a$, $B$--$b$, $C$--$c$, and $D$--$d$. Here $C$--$d$ is a dissatisfied couple, implying that $M$ is unstable. However, if $C$ and $d$ marry each other, and $c$ and $D$ marry each other, we obtain the stable marriage given by $A$--$a$, $B$--$b$, $C$--$d$, $D$--$c$.
\end{example}

The problem of finding stable marriages has several interesting applications, for example in matching medical graduates to residency positions in hospitals. It can be shown that for every choice of preference lists there exists at least one stable marriage. We will prove this by presenting an algorithm to find a stable marriage.

\begin{algorithm}[H]
\caption{The Proposal Algorithm (Gale--Shapley)}
\label{alg:proposal}
\begin{algorithmic}[1]
\Require Complete preference lists for $n$ men and $n$ women
\Ensure A stable marriage $M$
\State Start with all men and women free (unmarried)
\While{there exists a free man $X$ who has not proposed to every woman}
    \State $X$ proposes to the most desirable woman $x$ on his list who has not already rejected him
    \If{$x$ is free}
        \State $x$ accepts; $X$ and $x$ become engaged
    \ElsIf{$x$ prefers $X$ over her current fianc\'{e} $Y$}
        \State $x$ jilts $Y$ and becomes engaged to $X$; $Y$ becomes free
    \Else
        \State $x$ rejects $X$
    \EndIf
\EndWhile
\State The engaged pairs form a stable marriage
\end{algorithmic}
\end{algorithm}

The basic idea behind this algorithm can be summarized as ``man proposes, woman disposes'': each currently unattached man proposes to the most desirable woman on his list who has not already rejected him, and this woman then decides whether to accept or reject a proposal.

\begin{theorem}[Proposal Algorithm Terminates]
\label{thm:proposal-term}
The Proposal Algorithm terminates with a stable marriage in at most $n^2$ proposals.
\end{theorem}

\begin{proof}
A woman once married will stay married during the course of the algorithm, although her mates may change with time. Furthermore, the desirability of her mates (in her view) can only improve with time. Thus at each step either a woman gets married for the first time, or an already married woman obtains a more desirable mate.

An unattached man always has at least one woman available that he can proposition. This is because every woman he has already proposed to is currently married, and if he runs out of women then all women are married---this cannot happen unless all men are married too. Since at each step the proposer will eliminate one woman on his list, and the total size of the lists is $n^2$, we conclude that the algorithm uses at most $n^2$ proposals.

We claim that the final marriage $M$ is stable. Otherwise, let $X$--$y$ be a dissatisfied pair, where in $M$ they are paired as $X$--$x$ and $Y$--$y$. Since $X$ prefers $y$ to $x$, he must have proposed to $y$ before getting married to $x$. Since $y$ either rejected $X$, or accepted him only to jilt him later, her mates thereafter (including $Y$) must be more desirable to her than $X$. Therefore, $y$ must prefer $Y$ to $X$, contradicting the assumption that $y$ is dissatisfied.
\end{proof}

\subsection{Average-Case Analysis and the Principle of Deferred Decisions}

Our interest here is in performing an average-case analysis of this algorithm. Thus we are considering a probabilistic analysis of a deterministic algorithm. We introduce this analysis here because it touches upon several tools that are important in the analysis of randomized algorithms.

For this average-case analysis, we assume that the men's lists are chosen independently and uniformly at random; the women's lists can be arbitrary but must be fixed in advance. Let the random variable $T_P$ denote the number of proposals made during the execution of the Proposal Algorithm.

We present a very simple technique---the \emph{Principle of Deferred Decisions}---for getting around problems of dependency in the analysis. The idea is to not assume that the entire set of random choices is made in advance. Rather, at each step of the process we fix only the random choices that must be revealed to the algorithm.

The Principle of Deferred Decisions can be used to simplify the average-case analysis of the Proposal Algorithm as follows. We do not assume that the men have chosen their (random) preference list in advance. Each time a man has to make a proposal, he picks a random woman from the set of women not already propositioned by him, and proceeds to propose to her.

Suppose instead that each time a man makes a proposal, he chooses a woman uniformly at random from the set of all $n$ women, including those to whom he has already proposed. Call this new algorithm the \emph{Amnesiac Algorithm}. The output produced by the Amnesiac Algorithm is exactly the same as that of the original Proposal Algorithm. The only difference is that there are some wasted proposals in the Amnesiac Algorithm. Let $T_A$ denote the number of proposals made by the Amnesiac Algorithm. Clearly, $T_A$ stochastically dominates $T_P$: for all $m$, $\Pr[T_A > m] \geq \Pr[T_P > m]$.

The algorithm terminates with a stable marriage once all women have received at least one proposal each. As will become clear in the next section, bounding the value of $T_A$ is a special case of the Coupon Collector's Problem.

\begin{theorem}
\label{thm:stable-marr-prob}
For any constant $c \in \mathbb{R}$, and $m = n \ln n + cn$,
\[
\lim_{n \to \infty} \Pr[T_A > m] = 1 - e^{-e^{-c}}.
\]
\end{theorem}

\begin{cpplisting}
\label{code:stable-marriage}
\begin{lstlisting}[language=C++, frame=single, framerule=0.4pt,
  rulecolor=\color{gray}, backgroundcolor=\color{gray!5},
  basicstyle=\ttfamily\footnotesize, keywordstyle=\bfseries,
  commentstyle=\itshape\color{gray}, numbers=left,
  numberstyle=\tiny\color{gray}, numbersep=8pt,
  tabsize=2, showstringspaces=false]
#include <algorithm>
#include <iostream>
#include <random>
#include <vector>

struct StableMarriage {
    int n;
    std::vector<std::vector<int>> men_prefs;
    std::vector<std::vector<int>> women_prefs;
    std::vector<int> women_rank;  // rank[w][m]: how much w likes m
    std::vector<int> next_proposal;
    std::vector<int> current_wife;
    std::vector<int> current_husband;
    int proposals;

    StableMarriage(int n, std::mt19937& rng)
        : n(n), men_prefs(n), women_prefs(n),
          women_rank(n, std::vector<int>(n)),
          next_proposal(n, 0),
          current_wife(n, -1), current_husband(n, -1),
          proposals(0) {
        // Generate random preference lists
        for (int i = 0; i < n; ++i) {
            men_prefs[i].resize(n);
            for (int j = 0; j < n; ++j) men_prefs[i][j] = j;
            std::shuffle(men_prefs[i].begin(),
                         men_prefs[i].end(), rng);
            women_prefs[i].resize(n);
            for (int j = 0; j < n; ++j) women_prefs[i][j] = j;
            std::shuffle(women_prefs[i].begin(),
                         women_prefs[i].end(), rng);
            for (int j = 0; j < n; ++j)
                women_rank[i][women_prefs[i][j]] = j;
        }
    }

    void run() {
        std::vector<int> free_men(n);
        for (int i = 0; i < n; ++i) free_men[i] = i;
        while (!free_men.empty()) {
            int man = free_men.back();
            int woman = men_prefs[man][next_proposal[man]++];
            ++proposals;
            if (current_husband[woman] == -1) {
                current_husband[woman] = man;
                current_wife[man] = woman;
                free_men.pop_back();
            } else if (women_rank[woman][man]
                       < women_rank[woman][current_husband[woman]]) {
                int old = current_husband[woman];
                current_husband[woman] = man;
                current_wife[man] = woman;
                current_wife[old] = -1;
                free_men.pop_back();
                free_men.push_back(old);
            }
        }
    }

    bool is_stable() const {
        for (int m = 0; m < n; ++m) {
            int w = current_wife[m];
            for (int pref = 0; pref < women_rank[w][m]; ++pref) {
                int w_pref = women_prefs[w][pref];
                int m_curr = current_husband[w];
                if (women_rank[w_pref][w]
                    < women_rank[w_pref][current_wife[w_pref]])
                    return false;
            }
        }
        return true;
    }
};

int main() {
    std::mt19937 rng(42);
    std::cout << "Stable Marriage (Gale-Shapley):\n";
    for (int n : {4, 8, 16, 32, 64}) {
        double total_proposals = 0;
        int trials = 1000;
        for (int t = 0; t < trials; ++t) {
            StableMarriage sm(n, rng);
            sm.run();
            total_proposals += sm.proposals;
        }
        double avg = total_proposals / trials;
        double theory = n * std::log(n);
        std::cout << "  n=" << n
                  << "  avg_proposals=" << avg
                  << "  n*ln(n)=" << theory << "\n";
    }
    return 0;
}
\end{lstlisting}
\end{cpplisting}

\section{The Coupon Collector's Problem}
\label{sec:coupon}

In the coupon collector's problem, there are $n$ types of coupons and at each trial a coupon is chosen at random. Each random coupon is equally likely to be of any of the $n$ types, and the random choices of the coupons are mutually independent. Let $m$ be the number of trials. The goal is to study the relationship between $m$ and the probability of having collected at least one copy of each of the $n$ types. The reader may wish to make the correspondence between this process and an occupancy problem (Section~4.1) in which $m$ balls are randomly distributed in $n$ bins.

\subsection{An Elementary Analysis}

Let $X$ be a random variable defined to be the number of trials required to collect at least one of each type of coupon. We first determine the expected value of $X$.

Call the $i$th trial $C_i$ a \emph{success} if the type $C_i$ was not drawn in any of the first $i-1$ selections. We divide the sequence into \emph{epochs}, where epoch $i$ begins with the trial following the $i$th success and ends with the trial on which we obtain the $(i+1)$st success. Define the random variable $X_i$, for $0 \leq i \leq n-1$, to be the number of trials in the $i$th epoch, so that $X = \sum_{i=0}^{n-1} X_i$.

The random variable $X_i$ is geometrically distributed with parameter $p_i = (n-i)/n$ (the probability of drawing one of the $n-i$ remaining coupon types). Thus, the expected value of $X_i$ is $1/p_i = n/(n-i)$ and its variance is $(1-p_i)/p_i^2 = i/(n-i)$. By linearity of expectation,
\[
\mathbb{E}[X] = \mathbb{E}\Bigl[\sum_{i=0}^{n-1} X_i\Bigr] = \sum_{i=0}^{n-1} \mathbb{E}[X_i] = \sum_{i=1}^{n} \frac{n}{i} = n H_n.
\]
By Proposition~B.4 the $n$th Harmonic number $H_n$ is asymptotically equal to $\ln n + \Theta(1)$, implying that $\mathbb{E}[X] = n \ln n + O(n)$.

\begin{theorem}[Coupon Collector Expected Time]
\label{thm:coupon-expect}
The expected number of trials to collect all $n$ coupon types is $n H_n = n \ln n + O(n)$.
\end{theorem}

Since the $X_i$'s are independent, we can determine the variance of $X$ using the additivity of variances:
\[
\sigma_X^2 = \sum_{i=0}^{n-1} \sigma_{X_i}^2 = \sum_{i=0}^{n-1} \frac{i}{(n-i)^2} \leq \frac{n}{2} \sum_{i=1}^{n} \frac{1}{i^2} \leq \frac{\pi^2 n}{12}.
\]

This variance bound lets us apply Chebyshev's inequality to obtain a
weak but useful tail bound.  For $t = n \ln n$ (the approximate mean),
\[
\Pr\!\bigl[|X - n H_n| \geq t\bigr] \leq \frac{\sigma_X^2}{t^2}
= O\!\left(\frac{n}{n^2 \ln^2 n}\right) = O\!\left(\frac{1}{n \ln^2 n}\right).
\]
In particular, the probability that $X$ deviates from its mean by
more than $n \ln n$ decays like $O(1/(n \ln^2 n))$, which is much
weaker than the exponential bound obtained by the union bound below.
This illustrates a general principle: Chebyshev gives polynomial tail
bounds that are easy to derive, while more refined methods (Chernoff,
sharp threshold) give exponentially stronger guarantees.

\begin{exercise}
Use the Chebyshev inequality to find an upper bound on the probability that $X > \beta n \ln n$, for a constant $\beta > 1$.
\end{exercise}

Let $\mathcal{E}_i$ denote the event that coupon type $i$ is not collected in the first $r$ trials. Using the inequality $(1 - 1/n)^r \leq e^{-r/n}$, we obtain that
\[
\Pr[\mathcal{E}_i] = \Bigl(1 - \frac{1}{n}\Bigr)^r \leq e^{-r/n}.
\]
This bound is $n^{-\beta}$ for $r = \beta n \ln n$. Using the union bound, we obtain for $r = \beta n \ln n$,
\[
\Pr[X > r] = \Pr\Bigl[\bigcup_{i=1}^{n} \mathcal{E}_i\Bigr] \leq \sum_{i=1}^{n} \Pr[\mathcal{E}_i] \leq n \cdot n^{-\beta} = n^{1-\beta}.
\]

\subsection{The Poisson Heuristic}

Before we show the sharp concentration result for $X$, the following heuristic argument will help to establish some intuition. The heuristic argument is based on the approximation of the binomial distribution by the Poisson distribution.

Let $N_i^{(r)}$ denote the number of times the coupon of type $i$ is chosen during the first $r$ trials; the event $\mathcal{E}_i$ is the same as the event $\{N_i^{(r)} = 0\}$. The random variable $N_i^{(r)}$ has the binomial distribution with parameters $r$ and $p = 1/n$. For suitably small $\lambda$ and as $r$ approaches $\infty$, the Poisson distribution with parameter $\lambda$ is a good approximation to the binomial distribution with parameters $r$ and $p$. In the current setting, we approximate the distribution of $N_i^{(r)}$ by the Poisson distribution with parameter $\lambda = r/n$. A random variable $Y$ has the Poisson distribution with parameter $\lambda$ if for any non-negative integer $y$,
\[
\Pr[Y = y] = \frac{\lambda^y e^{-\lambda}}{y!}.
\]
Using this approximation,
\begin{equation}
\Pr[\mathcal{E}_i] = \Pr[N_i^{(r)} = 0] \approx \frac{\lambda^0 e^{-\lambda}}{0!} = e^{-r/n}.
\label{eq:poisson-approx}
\end{equation}

The main benefit of the Poisson approximation is that we can claim the events $\mathcal{E}_i$, for $1 \leq i \leq n$, are ``almost independent'' even though there is some dependence between them. If the approximation in~\eqref{eq:poisson-approx} were exact, we would obtain that the events $\mathcal{E}_i$ are truly independent. Making the heuristic assumption of independence, we obtain that for $1 \leq i \leq n$, the probability that all coupon types are collected in the first $m$ trials is given by
\[
\Pr\Bigl[\bigcap_{i=1}^{n} (\overline{\mathcal{E}_i})\Bigr] = \prod_{i=1}^{n} \Pr[\overline{\mathcal{E}_i}] \approx (1 - e^{-m/n})^n \approx e^{-n e^{-m/n}}.
\]
Let $m = n(\ln n + c)$ for any constant $c \in \mathbb{R}$. Then,
\[
\Pr[X > m] \approx \Pr\Bigl[\bigcup_{i=1}^{n} \mathcal{E}_i\Bigr] \approx 1 - e^{-e^{-c}}.
\]
Observe that this probability $e^{-e^{-c}}$ is close to $1$ for large positive $c$, and is negligibly small for large negative $c$. Thus, the probability of having collected all $n$ coupon types abruptly changes from nearly zero to almost one in a small interval centered around $n \ln n$.

\subsection{A Sharp Threshold}

We now convert the heuristic argument from the previous section into a rigorous (but significantly more complex) proof using the Boole--Bonferroni Inequalities (Proposition~C.2).

\begin{lemma}
\label{lem:coupon-tech}
Let $c$ be a real constant, and $m = n \ln n + cn$ for positive integer $n$. Then, for any fixed positive integer $k$,
\[
\lim_{n \to \infty} \binom{n}{k}\Bigl(1 - \frac{k}{n}\Bigr)^m = \frac{e^{-ck}}{k!}.
\]
\end{lemma}

\begin{theorem}[Sharp Threshold for Coupon Collector]
\label{thm:coupon-sharp}
Let the random variable $X$ denote the number of trials for collecting each of the $n$ types of coupons. Then, for any constant $c \in \mathbb{R}$, and $m = n \ln n + cn$,
\[
\lim_{n \to \infty} \Pr[X > m] = 1 - e^{-e^{-c}}.
\]
\end{theorem}

\begin{proof}
We have that the event $\{X > m\} = \bigcup_{i=1}^{n} \mathcal{E}_i$. By the Principle of Inclusion--Exclusion,
\[
\Pr\Bigl[\bigcup_{i=1}^{n} \mathcal{E}_i\Bigr] = \sum_{k=1}^{n} (-1)^{k+1} P_k,
\]
where $P_k = \binom{n}{k} (1 - k/n)^m$ is the $k$th elementary symmetric sum. Let $s_k = P_1 - P_2 + P_3 - \cdots + (-1)^{k+1} P_k$ denote the partial sum formed by the first $k$ terms of this series. By the Boole--Bonferroni inequalities (Proposition~C.2), we have the bracketing property of the partial sums:
\[
s_{2j} \leq \Pr\Bigl[\bigcup_{i=1}^{n} \mathcal{E}_i\Bigr] \leq s_{2j+1}.
\]
By Lemma~\ref{lem:coupon-tech}, for each $k$, $P_k \to e^{-ck}/k!$ as $n \to \infty$. Define the partial sums of the terms
\[
s_k = \sum_{j=1}^{k} (-1)^{j+1} \frac{e^{-ck}}{k!} = 1 - e^{-e^{-c}}.
\]
Notice that the right-hand side consists precisely of the first $k$ terms of the power series expansion of $f(c) = 1 - e^{-e^{-c}}$. We conclude that $\lim_{k \to \infty} s_k = f(c)$. The bracketing property of partial sums then implies that
\[
\Pr\Bigl[\bigcup_{i=1}^{n} \mathcal{E}_i\Bigr] \to f(c) = 1 - e^{-e^{-c}}
\]
as $n \to \infty$.
\end{proof}

By this theorem, for any real constant $c$, we have
\[
\lim_{n \to \infty} \Pr[X < n(\ln n - c)] = e^{-e^c}, \quad
\lim_{n \to \infty} \Pr[X > n(\ln n + c)] = 1 - e^{-e^{-c}}.
\]
Thus, with extremely high probability, the number of trials for collecting all $n$ coupon types lies in a small interval centered about its expected value. This result is almost like a deterministic result since it so sharply identifies the threshold value for collecting all coupons. We refer to such results as \emph{sharp threshold results}.

\begin{cpplisting}
\label{code:coupon}
\begin{lstlisting}[language=C++, frame=single, framerule=0.4pt,
  rulecolor=\color{gray}, backgroundcolor=\color{gray!5},
  basicstyle=\ttfamily\footnotesize, keywordstyle=\bfseries,
  commentstyle=\itshape\color{gray}, numbers=left,
  numberstyle=\tiny\color{gray}, numbersep=8pt,
  tabsize=2, showstringspaces=false]
#include <cmath>
#include <iostream>
#include <random>
#include <vector>

// Simulate the coupon collector: return the number of trials
// needed to collect all n coupon types.
int simulate_coupon_collector(int n, std::mt19937& rng) {
    std::uniform_int_distribution<int> dist(0, n - 1);
    std::vector<bool> collected(n, false);
    int remaining = n;
    int trials = 0;
    while (remaining > 0) {
        int coupon = dist(rng);
        ++trials;
        if (!collected[coupon]) {
            collected[coupon] = true;
            --remaining;
        }
    }
    return trials;
}

int main() {
    std::mt19937 rng(42);
    std::cout << "Coupon Collector Problem:\n";
    std::cout << "  n   avg_trials  n*H_n"
              << "  Pr[X>n(ln n+c)] theory vs sim\n";

    for (int n : {10, 50, 100, 500, 1000}) {
        double harmonic = 0.0;
        for (int i = 1; i <= n; ++i) harmonic += 1.0 / i;
        double expected_theory = n * harmonic;

        double total = 0;
        int trials_sim = 1000;
        for (int t = 0; t < trials_sim; ++t)
            total += simulate_coupon_collector(n, rng);
        double avg = total / trials_sim;

        // Check the sharp threshold for c=0 (m = n ln n)
        double m_threshold = n * std::log(n);
        int exceed_count = 0;
        for (int t = 0; t < trials_sim; ++t) {
            int tc = 0;
            std::vector<bool> seen(n, false);
            int rem = n;
            std::uniform_int_distribution<int> d(0, n - 1);
            while (rem > 0 && tc < static_cast<int>(m_threshold)) {
                int c = d(rng);
                ++tc;
                if (!seen[c]) { seen[c] = true; --rem; }
            }
            if (rem > 0) ++exceed_count;
        }
        double sim_prob = static_cast<double>(exceed_count) / trials_sim;
        double theory_prob = 1.0 - std::exp(-std::exp(0.0));
        // For c=0: limit = 1 - e^{-1} ~ 0.6321
        theory_prob = 1.0 - std::exp(-1.0);

        std::cout << "  " << n << "  " << avg
                  << "  " << expected_theory
                  << "  theory=" << theory_prob
                  << " sim=" << sim_prob << "\n";
    }
    return 0;
}
\end{lstlisting}
\end{cpplisting}

\addcontentsline{toc}{section}{Notes}
\section*{Notes}

Comprehensive treatises on occupancy problems are the books by Johnson and Kotz~\cite{JohnsonKotz1977}, and by Kolchin, Chistiakov, and Sevastianov~\cite{Kolchin1978}. Generally, we will be concerned with statements involving asymptotic estimates on random variables and probabilities. We will return to such estimates for occupancy problems in Chapter~5. Recent work by Azar, Broder, Karlin, and Upfal~\cite{Azar1999} builds on the basic occupancy problem and points out many applications to computer science.

The history of tail inequalities such as the Chebyshev bound dates back to the early days of probability theory. Following Chebyshev's bound~\cite{Chebyshev1867}, Markov~\cite{Markov1912} observed that the same idea could be used with higher moments. Kolmogorov~\cite{Kolmogorov1929} went further and remarked that $\Pr[X \geq r] \leq \mathbb{E}[f(X)]/s$ for any function $f(X)$, provided that $\mathbb{E}[f(X)]$ exists and $f(x) \geq s > 0$ for all $x \geq r$. The latter idea was exploited by Bernstein and by Chernoff in a manner we will describe in Chapter~5.

Classic sources for deterministic selection algorithms are the papers of Blum, Floyd, Pratt, Rivest, and Tarjan~\cite{Blum1973}, and of Sch\"{o}nhage, Paterson, and Pippenger~\cite{Schonhage1976}. The LazySelect algorithm presented here is a variant on one reported by Floyd and Rivest~\cite{Floyd1975}. The lower bound of $2n$ for median selection is due to Bent and John~\cite{Bent1985}.

The construction of pairwise independent random variables in Exercise~3.7 is given in Joffe~\cite{Joffe1946}. Its application to the reduction of random bits used by abstract randomized algorithms is due to Chor and Goldreich~\cite{ChorGoldreich1989}; Luby~\cite{Luby1988} presented this idea in the context of derandomizing algorithms.

The stable marriage problem was introduced by Gale and Shapley~\cite{GaleShapley1962}. The average-case analysis using the Principle of Deferred Decisions follows Knuth~\cite{Knuth1976}. The connection to the Coupon Collector's Problem was noted by Irving~\cite{Irving1985}.

The Coupon Collector's Problem is classical; see the books by Feller~\cite{Feller1968} and by Ross~\cite{Ross1996} for textbook treatments. The sharp threshold result of Theorem~3.8 is due to Erd\"{o}s and R\'{e}nyi~\cite{ErdosRenyi1961}.

\addcontentsline{toc}{section}{Problems}
\section*{Problems}

\begin{problem}
{\bf (Balls into Bins---Maximum Load)} Consider the experiment of throwing $m$ balls into $n$ bins, where each ball lands in a bin chosen independently and uniformly at random. Prove that if $m = n$, then the maximum load $L = \max_{1 \leq i \leq n} X_i$ satisfies
\[
\Pr\Bigl[L \geq \frac{3 \ln n}{\ln \ln n}\Bigr] \leq \frac{1}{n}.
\]
\end{problem}

\begin{problem}
{\bf (Birthday Paradox---Generalized)} Suppose that we have a universe of size $n$ and we sample $m$ elements independently and uniformly at random. Show that the probability of a collision (two sampled elements being equal) is at most $m^2 / (2n)$. Deduce that for $m = O(\sqrt{n})$, a collision is likely.
\end{problem}

\begin{problem}
{\bf (LazySelect---Expected Comparisons)} Prove rigorously that the expected number of comparisons performed by LazySelect is $2n + o(n)$. Your analysis should account for the expected number of retries.
\end{problem}

\begin{problem}
{\bf (Pairwise Independence)} Let $n$ be a prime and let $a, b$ be chosen independently and uniformly at random from $\mathbb{Z}_n$. For each $i \in \mathbb{Z}_n$, define $Y_i = ai + b \pmod{n}$. Prove that the $Y_i$'s are pairwise independent but not mutually independent.
\end{problem}

\begin{problem}
{\bf (Stable Marriage---Optimality)} Show that the Proposal Algorithm is \emph{male-optimal}: every man receives the best partner possible among all stable marriages. Show by example that it is not female-optimal.
\end{problem}

\begin{problem}
{\bf (Coupon Collector---Full Distribution)} Derive a more refined estimate of the distribution of $X$, the number of trials to collect all $n$ coupons. Specifically, show that for any constant $c \in \mathbb{R}$,
\[
\Pr[X \leq n \ln n + cn] \to e^{-e^{-c}} \quad \text{as } n \to \infty.
\]
\end{problem}

\begin{problem}
{\bf (Weighted Coupon Collector)} Suppose there are $n$ coupon types, but coupon $i$ is drawn with probability $p_i$ (where $\sum_i p_i = 1$). Let $X$ be the number of trials to collect all $n$ types. Show that
\[
\mathbb{E}[X] = \int_0^{\infty} \Bigl(1 - \prod_{i=1}^{n}(1 - e^{-p_i t})\Bigr)\, dt.
\]
\end{problem}

\begin{problem}
{\bf (Second Moment Method)} Let $X$ be a non-negative random variable with $\mu = \mathbb{E}[X] > 0$ and $\sigma^2 = \operatorname{Var}(X)$. Prove that
\[
\Pr[X > 0] \geq \frac{\mu^2}{\mu^2 + \sigma^2}.
\]
Show that this is tight.
\end{problem}

\begin{problem}
{\bf (Birthday Attack on Hash Functions)} Suppose $h : \{0,1\}^* \to \{0,1\}^k$ is a hash function that maps inputs to $n = 2^k$ values. We hash $m$ independent inputs. Using the birthday bound, show that we need $m \approx 2^{k/2}$ inputs to find a collision with probability $\approx 1/2$. What does this imply for the security of hash functions?
\end{problem}

\begin{problem}
{\bf (Amnesiac Algorithm)} Consider the Amnesiac Algorithm for the Stable Marriage Problem, where each man proposes to a uniformly random woman (ignoring past rejections). Prove formally that $T_A$ stochastically dominates $T_P$ (the number of proposals in the original Proposal Algorithm), i.e., $\Pr[T_A > m] \geq \Pr[T_P > m]$ for all $m$.
\end{problem}

\end{document}
